\documentclass[11pt, aspectratio=169]{beamer}

% Paquetes
\usepackage[utf8]{inputenc}
\usepackage[T1]{fontenc}
\usepackage[spanish]{babel}
\usepackage{ragged2e}  									% \justifying text
\usepackage{booktabs}  									% Tables
\usepackage{tabularx}
\usepackage{tikz}      										% Diagrams
\usetikzlibrary{calc, shapes, backgrounds}
\usepackage[american]{circuitikz}
\usepackage{amsmath, amssymb}
\usepackage{url}       											% \url`s
\usepackage{listings}  										% Code listings

\usepackage{dsfont}
\usepackage{algorithm, algpseudocode}
\usepackage{changepage}

%usco theme
\usetheme{usco}

% Portada
\course{Nombre del curso}
\title{Titulo de la presentación}
\date[\today]{\today}
\author[Nombre Apellido]{Nombre completo\inst{1}, {\small \textit{Eng. Ms.C, Ph.D}}}
\email{x@usco.edu.co}
\institute{\inst{1}Universidad Surcolombiana}
\summary{Resumen de la presentación}

\begin{document}

\begin{frame}
\titlepage
\end{frame}

% ------------------------------------------------- CONTENIDO  -------------------------------------------------
\begin{frame}{Contenido}
    \tableofcontents
\end{frame}

% ------------------------------------------------- INTRODUCCIÓN  -------------------------------------------------
\section{Listas y enumeraciones}
\begin{frame}{Listas y enumeraciones}
    \begin{columns}[onlytextwidth]
    \column{.5\textwidth}
		\begin{itemize}
		    \item Contexto general del trabajo
		    \begin{itemize}
		        \item Descripción del tema principal.
		        \item Importancia del problema abordado.
		        \item Relación con el área de estudio.
		    \end{itemize}
		
		    \item Motivación
		    \begin{itemize}
		        \item Razones que justifican el desarrollo del trabajo.
		        \item Necesidad académica, técnica o institucional.
		    \end{itemize}
		
		    \item Alcance
		    \begin{itemize}
		        \item Aspectos que serán tratados en la presentación.
		        \item Límites del trabajo desarrollado.
		    \end{itemize}
		\end{itemize}
	\column{.5\textwidth}
		\begin{enumerate}
		    \item Contexto general del trabajo
		    \begin{enumerate}
		        \item Descripción del tema principal.
		        \item Importancia del problema abordado.
		        \item Relación con el área de estudio.
		    \end{enumerate}
		
		    \item Motivación
		    \begin{enumerate}
		        \item Razones que justifican el desarrollo del trabajo.
		        \item Necesidad académica, técnica o institucional.
		    \end{enumerate}
		
		    \item Alcance
		    \begin{enumerate}
		        \item Aspectos que serán tratados en la presentación.
		        \item Límites del trabajo desarrollado.
		    \end{enumerate}
		\end{enumerate}
    \end{columns}
\end{frame}


% ------------------------------------------------- OBJETIVOS  -------------------------------------------------
\section{Elementos estructurales}

\subsection{Bloques de texto}
\begin{frame}{Bloques de texto}{Tipo: simple, ejemplo, y \alert{alerta}}
    \begin{columns}[onlytextwidth]
	    \column{0.48\textwidth}
	    \begin{block}{Un bloque simple (block)}\small
	    		Este es un bloque de texto simple (\alert{block}).
	    \end{block}
	    \begin{exampleblock}{Un bloque de ejemplo (exampleblock)}\small
	    		Este es un bloque de texto del tipo ejemplo (\alert{exampleblock}).
		\end{exampleblock}
	    \begin{alertblock}{Un bloque de alerta (alertblock)}\small
	    		Este es un bloque de texto del tipo alerta (\alert{alertblock}).
	    \end{alertblock}
	\column{.48\textwidth}
	    \begin{block}{}\small
	    		Este es un bloque de texto simple (\alert{block}).
	    \end{block}
	    \begin{exampleblock}{}\small
	    		Este es un bloque de texto del tipo ejemplo (\alert{exampleblock}).
		\end{exampleblock}
	    \begin{alertblock}{}\small
	    		Este es un bloque de texto del tipo alerta (\alert{alertblock}).
	    \end{alertblock}
	 \end{columns}
\end{frame}

\begin{frame}{Otros bloques de texto}
    \begin{columns}[onlytextwidth]
	    \column{0.67\textwidth}
	    \begin{brownblock}{Un bloque brownblock}\small
	    		Este es un bloque de texto del tipo brownblock (\alert{brownblock}).
	    \end{brownblock}
	    \begin{darkredblock}{Un bloque darkredblock}\small
	    		Este es un bloque de texto del tipo darkredblock (\alert{darkredblock}).
		\end{darkredblock}
	    \begin{lightredblock}{Un bloque lightredblock}\small
	    		Este es un bloque de texto del tipo lightredblock (\alert{lightredblock}).
	    \end{lightredblock}
	    \begin{lightgreenblock}{Un bloque lightgreenblock}\small
	    		Este es un bloque de texto del tipo lightgreenblock (\alert{lightgreenblock}).
	    \end{lightgreenblock}
	\column{.28\textwidth}
	    \begin{brownblock}{}\tiny
	    		Este es un bloque de texto del tipo brownblock (\alert{brownblock}).
	    \end{brownblock}
	    \begin{darkredblock}{}\tiny
	    		Este es un bloque de texto del tipo darkredblock (\alert{darkredblock}).
		\end{darkredblock}
	    \begin{lightredblock}{}\tiny
	    		Este es un bloque de texto del tipo lightredblock (\alert{lightredblock}).
	    \end{lightredblock}
	    \begin{lightgreenblock}{}\tiny
	    		Este es un bloque de texto del tipo lightgreenblock (\alert{lightgreenblock}).
	    \end{lightgreenblock}
	\end{columns}
\end{frame}

\begin{frame}{Otros bloques de texto}
    \begin{rededgeblock}{rededgeblock}
	    		Este es un bloque de texto del tipo rededgeblock (\alert{rededgeblock}).
    \end{rededgeblock}
    
    \begin{rededgeblock_sm}{rededgeblock\_sm}
	    		Este es un bloque de texto del tipo rededgeblock\_sm (\alert{rededgeblock\_sm}).
    \end{rededgeblock_sm}
\end{frame}


% ------------------------------------------------- Definiciones, teoremas y demostraciones  -------------------------------------------------
\section{Definiciones, teoremas y demostraciones}

\begin{frame}{Definitions, theorems, and proofs}   
	\begin{definition}
		$\forall a,b\in\mathds{Z}: a\mid b\iff\exists c\in\mathds{Z}:a\cdot c=b$
	\end{definition}
	\begin{theorem}
		$\forall a\in\mathds{Z}: a\mid 0$
	\end{theorem}
	\begin{proof}
		$\forall a\in\mathds{Z}: a\cdot 0=0$
	\end{proof}
\end{frame}

% ------------------------------------------------- Números and matemáticas  -------------------------------------------------
\subsection{Números and matemáticas}
\begin{frame}[label=math]{Números and matemáticas}
    \framesubtitle{Formulas, ecuaciones, y expresiones}
    \begin{columns}[onlytextwidth]
    \column{.25\textwidth}
        1234567890
    \column{.25\textwidth}
        \oldstylenums{1234567890}
    \column{.25\textwidth}
        $\hat{x}$, $\check{x}$, $\tilde{a}$,
        $\bar{a}$, $\dot{y}$, $\ddot{y}$
    \column{.45\textwidth}
        $\int \!\! \int f(x,y,z)\,\mathsf{d}x\mathsf{d}y\mathsf{d}z$
    \end{columns}
    \vspace{1em}
    \begin{columns}[onlytextwidth]
    \column{.5\textwidth}
        $$\frac{1}{\displaystyle 1+
        \frac{1}{\displaystyle 2+
        \frac{1}{\displaystyle 3+x}}} +
        \frac{1}{1+\frac{1}{2+\frac{1}{3+x}}}$$
    \column{.5\textwidth}
        $$F:\left| \begin{array}{ccc}
        F''_{xx} & F''_{xy} &  F'_x \\
        F''_{yx} & F''_{yy} &  F'_y \\
        F'_x     & F'_y     & 0
        \end{array}\right| = 0$$
    \end{columns}
    \vspace{1em}
    \begin{columns}[onlytextwidth]
    \column{.33\textwidth}
        $$\mathop{\int \!\!\! \int}_{\mathbf{x} \in \mathds{R}^2}
        \! \langle \mathbf{x},\mathbf{y}\rangle\,\mathsf{d}\mathbf{x}$$
    \column{.33\textwidth}
        $$\overline{\overline{a\alpha}^2+\underline{b\beta}
        +\overline{\overline{d\delta}}}$$
    \column{.33\textwidth}
        $\left] 0,1\right[ + \lceil x \rfloor - \langle x,y\rangle$
    \end{columns}
    \vspace{1em}
    \begin{columns}[onlytextwidth]
    \column{.5\textwidth}
        \begin{eqnarray*}
        e^x &\approx& 1+x+x^2/2! + \\
            && {}+x^3/3! + x^4/4!
        \end{eqnarray*}
    \column{.5\textwidth}
        $${n+1\choose k} = {n\choose k} + {n \choose k-1}$$
    \end{columns}
\end{frame}


\section{Figuras y código}
\subsection{Tablas}
\begin{frame}{Tablas}
    \begin{table}[!b]
    \begin{tabularx}{\textwidth}{Xrrr}
        \textbf{Columna 1} & \textbf{Columna 2} & \textbf{Columna 3} & \textbf{Columna 4} \\
        \toprule
        Uno       	& 1234  & 4567   & 59.09 \\
        Dos 			& 123    & 5678   & 14.90 \\
        Tres 			& 34      & 2345   & 1.39  \\
        Cuatro 		& 32      & 12444 & 0.69  \\
        \bottomrule
    \end{tabularx}
    \caption{Datos numéricos ultimo año.}
    \end{table}
\end{frame}

\subsection{Figuras}
\begin{frame}{Figuras}
    \begin{figure}[b]
    \centering
    \begin{circuitikz}[american, thick]
		%-----------------------------------------
		% Coordenadas principales
		%-----------------------------------------
		\coordinate (Ei) at (0,0);
		\coordinate (A)  at (0.75,0);     % nodo de entrada
		\coordinate (B)  at (3.4,0);     % nodo final red de entrada
		\coordinate (N)  at (4.25,0);    % nodo inversor / realimentación
		
		% Altura de R1-C1 y R4
		\coordinate (Aup) at (0.75,1.35);
		\coordinate (Bup) at (3.4,1.35);
		\coordinate (Nmid) at (4.25,1.35);
		
		% Altura de R2-C2
		\coordinate (Ntop) at (4.25,2.35);
		
		%-----------------------------------------
		% Amplificador operacional
		%-----------------------------------------
		\node[op amp, anchor=-, scale=1.0] (opamp) at (N) {};
		
		% Nodo de salida
		\coordinate (O) at ($(opamp.out)+(0.8,0)$);
		\coordinate (Omid) at (O |- Nmid);
		\coordinate (Otop) at (O |- Ntop);
		
		%-----------------------------------------
		% Entrada
		%-----------------------------------------
		\draw (Ei) node[ocirc]{} node[left]{\LARGE $e_i$}
		      -- (A) node[circ]{};
		
		%-----------------------------------------
		% Red de entrada: R3 en paralelo con R1-C1
		%-----------------------------------------
		\draw (A) to[R,l_=$R_3$] (B) node[circ]{}
		      -- (N) node[circ]{};
		
		\draw (A) -- (Aup)
		      to[R,l=$R_1$] ++(1.75,0)
		      to[C,l=$C_1$] (Bup)
		      -- (B);
		
		%-----------------------------------------
		% Red de realimentación: R4 en paralelo con R2-C2
		%-----------------------------------------
		\draw (N) -- (Ntop);
		
		\draw (Ntop)
		      to[R,l=$R_2$] ++(1.75,0)
		      to[C,l=$C_2$] (Otop)
		      -- (O);
		
		\draw (Nmid)
		      to[R,l_=$R_4$] (Omid);
		
		\draw (O) -- (Otop);
		
		%-----------------------------------------
		% Salida
		%-----------------------------------------
		\draw (opamp.out) -- (O) node[circ]{}
		      -- ++(0.65,0) node[ocirc]{} node[right]{\LARGE $e_o$};
		
		%-----------------------------------------
		% Entrada no inversora a tierra
		%-----------------------------------------
		\coordinate (G) at ($(opamp.+)+(0,-1.15)$);
		\draw (opamp.+) -- (G);
		
		%-----------------------------------------
		% Línea inferior de referencia
		%-----------------------------------------
		\coordinate (BotL) at (0,-1.8);
		\coordinate (BotR) at ($(O)+(0.65,-1.3)$);
		\coordinate (GB) at (G |- BotL);
		
		\draw (BotL) node[ocirc]{}
		      -- (GB) node[circ]{}
		      -- (BotR) node[ocirc]{};
		
		\draw (G) -- (GB);
		\draw (GB) -- ++(0,-0.15) node[ground]{};
		
		\end{circuitikz}
    \caption{Circuito con amplificadores operacionales.}
    \end{figure}
\end{frame}

\subsection{Código}
% Colores suaves para código
\definecolor{codebg}{RGB}{248,248,248}
\definecolor{codeframe}{RGB}{220,220,220}
\definecolor{codekeyword}{RGB}{0,82,155}
\definecolor{codestring}{RGB}{170,45,45}
\definecolor{codecomment}{RGB}{90,120,90}
\definecolor{codenumber}{RGB}{120,120,120}

\lstdefinestyle{pythonstyle}{
    language=Python,
    backgroundcolor=\color{codebg},
    basicstyle=\ttfamily\footnotesize,
    keywordstyle=\color{codekeyword}\bfseries,
    stringstyle=\color{codestring},
    commentstyle=\color{codecomment}\itshape,
    numberstyle=\tiny\color{codenumber},
    numbers=left,
    numbersep=8pt,
    frame=single,
    rulecolor=\color{codeframe},
    framerule=0.4pt,
    breaklines=true,
    showstringspaces=false,
    tabsize=4,
    xleftmargin=1.2em,
    framexleftmargin=1.2em,
    aboveskip=0.8em,
    belowskip=0.8em
}

\begin{frame}[fragile]{Código}{Ejemplo en Python: par o impar}
\begin{lstlisting}[style=pythonstyle]
for i in range(11):
    if i % 2 == 0:
        print(i, "es par")
    else:
        print(i, "es impar")
\end{lstlisting}
\end{frame}

\subsection{Algoritmos}
\begin{frame}{Gradiente descendiente}
	\begin{algorithm}[H]
		\caption{Algoritmo de descenso más pronunciado para una red neuronal}
		\label{alg:steepest-descent}

		\begin{algorithmic}[1]
			\State Inicializar pesos $W$ y sesgos $b$
			\State Definir la tasa de aprendizaje $\alpha$
			\State Definir el número máximo de iteraciones $N$
			
			\For{$k = 1,2,\dots,N$}
			    \State Calcular la salida de la red: $\hat{y}=f(x;W,b)$
			    \State Calcular la función de pérdida: $J(W,b)$
			    \State Calcular los gradientes: $\nabla_W J$ y $\nabla_b J$
			    \State Actualizar los pesos: $W \leftarrow W - \alpha \nabla_W J$
			    \State Actualizar los sesgos: $b \leftarrow b - \alpha \nabla_b J$
			\EndFor
			
			\State Retornar los parámetros entrenados $W$ y $b$
		\end{algorithmic}
	\end{algorithm}
\end{frame}

\section{Citaciones bibliográfica}

\begin{frame}{Citas y referencias} 
El método de gradiente descendente es un algoritmo iterativo usado para entrenar redes neuronales mediante la minimización de una función de pérdida (costo). Para ello, se actualizan los pesos en dirección opuesta al gradiente, reduciendo progresivamente el error del modelo \cite{rumelhart1986}. Sus variantes estocásticas, fundamentadas en la aproximación estocástica \cite{robbins1951}, permiten entrenar modelos a gran escala \cite{bottou2010}, mientras que métodos como \textit{Momentum}, \textit{RMSProp} y \textit{Adam} mejoran la convergencia y estabilidad del entrenamiento \cite{ruder2016}.
\end{frame}

\begin{frame}{Referencias}{\TeX, \LaTeX, and Beamer}
	\begin{thebibliography}{9} 
		\bibitem{robbins1951} H. Robbins and S. Monro, ``A Stochastic Approximation Method,'' \textit{The Annals of Mathematical Statistics}, vol. 22, no. 3, pp. 400--407, 1951.
		\bibitem{rumelhart1986} D. E. Rumelhart, G. E. Hinton, and R. J. Williams, ``Learning representations by back-propagating errors,'' \textit{Nature}, vol. 323, pp. 533--536, 1986. 
		\bibitem{bottou2010} L. Bottou, ``Large-Scale Machine Learning with Stochastic Gradient Descent,'' in \textit{Proceedings of COMPSTAT 2010}, pp. 177--186, 2010. 
		\bibitem{ruder2016} S. Ruder, ``An overview of gradient descent optimization algorithms,'' \textit{arXiv preprint arXiv:1609.04747}, 2016. \end{thebibliography} 
\end{frame}

% ------------------------------------------------- AGRADECIMIENTO  -------------------------------------------------
\agradecimiento

\end{document}


